
\documentclass{article}
\usepackage{colortbl}
\usepackage{makecell}
\usepackage{multirow}
\usepackage{supertabular}

\begin{document}

\newcounter{utterance}

\centering \large Interaction Transcript for game `codenames', experiment `opponent\_difficult', episode 2 with Qwen3.6{-}35B{-}A3B{-}FP8{-}without{-}reasoning.
\vspace{24pt}

{ \footnotesize  \setcounter{utterance}{1}
\setlength{\tabcolsep}{0pt}
\begin{supertabular}{c@{$\;$}|p{.15\linewidth}@{}p{.15\linewidth}p{.15\linewidth}p{.15\linewidth}p{.15\linewidth}p{.15\linewidth}}
    \# & $\;$A & \multicolumn{4}{c}{Game Master} & $\;\:$B\\
    \hline

    \theutterance \stepcounter{utterance}  
    & & & \multicolumn{2}{p{0.3\linewidth}}{
        \cellcolor[rgb]{0.9,0.9,0.9}{
            \makecell[{{p{\linewidth}}}]{
                \texttt{\tiny{[GM$|$GM]}}
                {'hidden': {'team': ['струна', 'тарілка', 'страйк', 'лід', 'єгипет', 'боєць', 'вашингтон', 'вчитель', 'храм'], 'innocent': ['весна', 'тінь', 'готель', 'заряд', 'костюм', 'орел', 'ліс'], 'opponent': ['європа', 'агент', 'театр', 'оливка', 'токіо', 'центр', "м'ята", 'бермуди'], 'assassin': ['бомба']}, 'revealed': {'team': {'team': [], 'innocent': [], 'opponent': [], 'assassin': []}, 'opponent': {'team': [], 'innocent': [], 'opponent': [], 'assassin': []}}}
            }
        }
    }
    & & \\ \\

    \theutterance \stepcounter{utterance}  
    & & \multicolumn{4}{p{0.6\linewidth}}{
        \cellcolor[rgb]{0.9,0.9,0.9}{
            \makecell[{{p{\linewidth}}}]{
                \texttt{\tiny{[P1$\langle$GM]}}
                \texttt{Надайте одну однословну підказку, яка пов'язана з деякими словами вашої команди зі списку нижче. Ви можете обрати, зі скількома і з якими саме словами пов'язана ваша підказка.} \\
\texttt{Мета полягає в тому, щоб ваш партнер відгадав усі слова вашої команди на основі підказок, які ви йому даєте, раніше, ніж інша команда відкриє всі свої. Якщо ваш партнер відгадає слово «суперника» або «відволікача», його хід відгадування закінчується, і команда суперників отримує право зробити хід. Якщо ваш партнер відгадає слово «вбивці», гра закінчується, і ваша команда негайно програє.} \\
\texttt{Чим із більшою кількістю слів пов'язана ваша підказка, тим швидше ви можете виграти гру, але тим складнішими можуть бути асоціації для відгадування вашим партнером, тож обирайте слово{-}підказку мудро. Слово{-}підказка має бути семантично пов'язаним із цільовими словами; воно не може бути одним зі слів у списках або містити їхні частини.} \\
\\ 
\texttt{Завжди надавайте вашу однословну підказку та ваш розділений комами список пов'язаних цільових слів у такому форматі й робіть відповіді якомога коротшими; ніколи не додавайте жодного іншого тексту, окрім того, що потрібно в цій формі:} \\
\\ 
\texttt{ПІДКАЗКА: <СЛОВО>} \\
\texttt{ЦІЛІ: <ЦІЛІ>} \\
\\ 
\texttt{Слова вашої команди: струна, тарілка, страйк, лід, єгипет, боєць, вашингтон, вчитель, храм.} \\
\texttt{Слова команди суперника: європа, агент, театр, оливка, токіо, центр, м'ята, бермуди.} \\
\texttt{Слова{-}відволікачі: весна, тінь, готель, заряд, костюм, орел, ліс.} \\
\texttt{Слова{-}вбивці: бомба.} \\
            }
        }
    }
    & \\ \\

    \theutterance \stepcounter{utterance}  
    & \multicolumn{4}{p{0.6\linewidth}}{
        \cellcolor[rgb]{0.9,0.9,0.9}{
            \makecell[{{p{\linewidth}}}]{
                \texttt{\tiny{[P1$\rangle$GM]}}
                \texttt{ПІДКАЗКА: музика} \\
\texttt{ЦІЛІ: струна, тарілка} \\
            }
        }
    }
    & & \\ \\

    \theutterance \stepcounter{utterance}  
    & & & \multicolumn{2}{p{0.3\linewidth}}{
        \cellcolor[rgb]{0.9,0.9,0.9}{
            \makecell[{{p{\linewidth}}}]{
                \texttt{\tiny{[GM$|$GM]}}
                \texttt{музика} \\
            }
        }
    }
    & & \\ \\

    \theutterance \stepcounter{utterance}  
    & & & \multicolumn{2}{p{0.3\linewidth}}{
        \cellcolor[rgb]{0.9,0.9,0.9}{
            \makecell[{{p{\linewidth}}}]{
                \texttt{\tiny{[GM$|$GM]}}
                ['струна', 'тарілка']
            }
        }
    }
    & & \\ \\

    \theutterance \stepcounter{utterance}  
    & & & \multicolumn{2}{p{0.3\linewidth}}{
        \cellcolor[rgb]{0.9,0.9,0.9}{
            \makecell[{{p{\linewidth}}}]{
                \texttt{\tiny{[GM$|$GM]}}
                {'word': 'струна', 'assignment': 'team'}
            }
        }
    }
    & & \\ \\

    \theutterance \stepcounter{utterance}  
    & & & \multicolumn{2}{p{0.3\linewidth}}{
        \cellcolor[rgb]{0.9,0.9,0.9}{
            \makecell[{{p{\linewidth}}}]{
                \texttt{\tiny{[GM$|$GM]}}
                {'word': 'тарілка', 'assignment': 'team'}
            }
        }
    }
    & & \\ \\

    \theutterance \stepcounter{utterance}  
    & & \multicolumn{4}{p{0.6\linewidth}}{
        \cellcolor[rgb]{0.9,0.9,0.9}{
            \makecell[{{p{\linewidth}}}]{
                \texttt{\tiny{[GM$\rangle$P2]}}
                \texttt{Надайте розділений комами список із не більш ніж 2 слів із наведеного нижче списку, які найкраще пов'язані або найбільш тісно асоціюються зі словом «музика». Завжди починайте свій список здогадок із ВІДГАДКА:  і не додавайте жодного іншого тексту у своїй відповіді.} \\
\\ 
\texttt{заряд, готель, храм, страйк, боєць, м'ята, бомба, весна, орел, костюм, лід, оливка, ліс, єгипет, бермуди, тарілка, вашингтон, європа, струна, вчитель, центр, тінь, токіо, агент, театр} \\
            }
        }
    }
    & \\ \\

    \theutterance \stepcounter{utterance}  
    & & & \multicolumn{4}{p{0.6\linewidth}}{
        \cellcolor[rgb]{0.9,0.9,0.9}{
            \makecell[{{p{\linewidth}}}]{
                \texttt{\tiny{[GM$\langle$P2]}}
                \texttt{ВІДГАДКА: струна, театр} \\
            }
        }
    }
     \\ \\

    \theutterance \stepcounter{utterance}  
    & & & \multicolumn{2}{p{0.3\linewidth}}{
        \cellcolor[rgb]{0.9,0.9,0.9}{
            \makecell[{{p{\linewidth}}}]{
                \texttt{\tiny{[GM$|$GM]}}
                ['струна', 'театр']
            }
        }
    }
    & & \\ \\

    \theutterance \stepcounter{utterance}  
    & & & \multicolumn{2}{p{0.3\linewidth}}{
        \cellcolor[rgb]{0.9,0.9,0.9}{
            \makecell[{{p{\linewidth}}}]{
                \texttt{\tiny{[GM$|$GM]}}
                {'word': 'струна', 'assignment': 'team'}
            }
        }
    }
    & & \\ \\

    \theutterance \stepcounter{utterance}  
    & & & \multicolumn{2}{p{0.3\linewidth}}{
        \cellcolor[rgb]{0.9,0.9,0.9}{
            \makecell[{{p{\linewidth}}}]{
                \texttt{\tiny{[GM$|$GM]}}
                {'word': 'струна', 'assignment': 'team'}
            }
        }
    }
    & & \\ \\

    \theutterance \stepcounter{utterance}  
    & & & \multicolumn{2}{p{0.3\linewidth}}{
        \cellcolor[rgb]{0.9,0.9,0.9}{
            \makecell[{{p{\linewidth}}}]{
                \texttt{\tiny{[GM$|$GM]}}
                {'word': 'театр', 'assignment': 'opponent'}
            }
        }
    }
    & & \\ \\

    \theutterance \stepcounter{utterance}  
    & & & \multicolumn{2}{p{0.3\linewidth}}{
        \cellcolor[rgb]{0.9,0.9,0.9}{
            \makecell[{{p{\linewidth}}}]{
                \texttt{\tiny{[GM$|$GM]}}
                \texttt{театр} \\
            }
        }
    }
    & & \\ \\

    \theutterance \stepcounter{utterance}  
    & & & \multicolumn{2}{p{0.3\linewidth}}{
        \cellcolor[rgb]{0.9,0.9,0.9}{
            \makecell[{{p{\linewidth}}}]{
                \texttt{\tiny{[GM$|$GM]}}
                {'word': 'бермуди', 'assignment': 'opponent'}
            }
        }
    }
    & & \\ \\

    \theutterance \stepcounter{utterance}  
    & & & \multicolumn{2}{p{0.3\linewidth}}{
        \cellcolor[rgb]{0.9,0.9,0.9}{
            \makecell[{{p{\linewidth}}}]{
                \texttt{\tiny{[GM$|$GM]}}
                {'word': 'європа', 'assignment': 'opponent'}
            }
        }
    }
    & & \\ \\

    \theutterance \stepcounter{utterance}  
    & & & \multicolumn{2}{p{0.3\linewidth}}{
        \cellcolor[rgb]{0.9,0.9,0.9}{
            \makecell[{{p{\linewidth}}}]{
                \texttt{\tiny{[GM$|$GM]}}
                {'hidden': {'team': ['тарілка', 'страйк', 'лід', 'єгипет', 'боєць', 'вашингтон', 'вчитель', 'храм'], 'innocent': ['весна', 'тінь', 'готель', 'заряд', 'костюм', 'орел', 'ліс'], 'opponent': ['агент', 'оливка', 'токіо', 'центр', "м'ята"], 'assassin': ['бомба']}, 'revealed': {'team': {'team': ['струна'], 'innocent': [], 'opponent': ['театр'], 'assassin': []}, 'opponent': {'team': [], 'innocent': [], 'opponent': ['бермуди', 'європа'], 'assassin': []}}}
            }
        }
    }
    & & \\ \\

    \theutterance \stepcounter{utterance}  
    & & \multicolumn{4}{p{0.6\linewidth}}{
        \cellcolor[rgb]{0.9,0.9,0.9}{
            \makecell[{{p{\linewidth}}}]{
                \texttt{\tiny{[P1$\langle$GM]}}
                \texttt{Слово струна було відгадано правильно. Слово театр було відгадано, але це суперницьке слово. Хід вашого товариша по команді закінчився.} \\
\texttt{Тепер надайте ще одну підказку, пов'язану з деякими з ваших командних слів, що залишилися, та список пов'язаних цільових слів. Пам'ятайте: почніть підказку з 'ПІДКАЗКА: ', перейдіть на новий рядок і почніть ваш список цільових слів, розділених комами, з 'ЦІЛІ: '. Зверніть увагу: деякі слова було вилучено зі списків порівняно з попередніми запитами.} \\
\\ 
\texttt{Ваші командні слова, що залишилися: тарілка, страйк, лід, єгипет, боєць, вашингтон, вчитель, храм.} \\
\texttt{Слова, що залишилися для вашого суперника: європа, агент, оливка, токіо, центр, м'ята, бермуди.} \\
\texttt{Слова{-}відволікачі, що залишилися: весна, тінь, готель, заряд, костюм, орел, ліс.} \\
\texttt{Слова{-}вбивці, що залишилися: бомба.} \\
            }
        }
    }
    & \\ \\

    \theutterance \stepcounter{utterance}  
    & \multicolumn{4}{p{0.6\linewidth}}{
        \cellcolor[rgb]{0.9,0.9,0.9}{
            \makecell[{{p{\linewidth}}}]{
                \texttt{\tiny{[P1$\rangle$GM]}}
                \texttt{ПІДКАЗКА: бойовий} \\
\texttt{ЦІЛІ: боєць, страйк} \\
            }
        }
    }
    & & \\ \\

    \theutterance \stepcounter{utterance}  
    & & & \multicolumn{2}{p{0.3\linewidth}}{
        \cellcolor[rgb]{0.9,0.9,0.9}{
            \makecell[{{p{\linewidth}}}]{
                \texttt{\tiny{[GM$|$GM]}}
                \texttt{бойовий} \\
            }
        }
    }
    & & \\ \\

    \theutterance \stepcounter{utterance}  
    & & & \multicolumn{2}{p{0.3\linewidth}}{
        \cellcolor[rgb]{0.9,0.9,0.9}{
            \makecell[{{p{\linewidth}}}]{
                \texttt{\tiny{[GM$|$GM]}}
                ['боєць', 'страйк']
            }
        }
    }
    & & \\ \\

    \theutterance \stepcounter{utterance}  
    & & & \multicolumn{2}{p{0.3\linewidth}}{
        \cellcolor[rgb]{0.9,0.9,0.9}{
            \makecell[{{p{\linewidth}}}]{
                \texttt{\tiny{[GM$|$GM]}}
                {'word': 'боєць', 'assignment': 'team'}
            }
        }
    }
    & & \\ \\

    \theutterance \stepcounter{utterance}  
    & & & \multicolumn{2}{p{0.3\linewidth}}{
        \cellcolor[rgb]{0.9,0.9,0.9}{
            \makecell[{{p{\linewidth}}}]{
                \texttt{\tiny{[GM$|$GM]}}
                {'word': 'страйк', 'assignment': 'team'}
            }
        }
    }
    & & \\ \\

    \theutterance \stepcounter{utterance}  
    & & \multicolumn{4}{p{0.6\linewidth}}{
        \cellcolor[rgb]{0.9,0.9,0.9}{
            \makecell[{{p{\linewidth}}}]{
                \texttt{\tiny{[GM$\rangle$P2]}}
                \texttt{Слово струна було відгадано правильно. Слово театр було відгадано, але це суперницьке слово. Ваш хід закінчився.} \\
\texttt{Тепер надайте ще один список слів, розділених комами, щонайменше з 1 і щонайбільше з 2 слів із наведеного нижче списку слів, які найкраще пов'язані або найбільш тісно асоціюються зі словом 'бойовий'. Пам'ятайте, що потрібно почати вашу відповідь з ВІДГАДКА: . Зауваження: деякі слова було вилучено зі списку порівняно з попередніми запитами.} \\
\\ 
\texttt{заряд, готель, храм, страйк, боєць, м'ята, бомба, весна, орел, костюм, лід, оливка, ліс, єгипет, тарілка, вашингтон, вчитель, центр, тінь, токіо, агент} \\
            }
        }
    }
    & \\ \\

    \theutterance \stepcounter{utterance}  
    & & & \multicolumn{4}{p{0.6\linewidth}}{
        \cellcolor[rgb]{0.9,0.9,0.9}{
            \makecell[{{p{\linewidth}}}]{
                \texttt{\tiny{[GM$\langle$P2]}}
                \texttt{ВІДГАДКА: боєць, заряд} \\
            }
        }
    }
     \\ \\

    \theutterance \stepcounter{utterance}  
    & & & \multicolumn{2}{p{0.3\linewidth}}{
        \cellcolor[rgb]{0.9,0.9,0.9}{
            \makecell[{{p{\linewidth}}}]{
                \texttt{\tiny{[GM$|$GM]}}
                ['боєць', 'заряд']
            }
        }
    }
    & & \\ \\

    \theutterance \stepcounter{utterance}  
    & & & \multicolumn{2}{p{0.3\linewidth}}{
        \cellcolor[rgb]{0.9,0.9,0.9}{
            \makecell[{{p{\linewidth}}}]{
                \texttt{\tiny{[GM$|$GM]}}
                {'word': 'боєць', 'assignment': 'team'}
            }
        }
    }
    & & \\ \\

    \theutterance \stepcounter{utterance}  
    & & & \multicolumn{2}{p{0.3\linewidth}}{
        \cellcolor[rgb]{0.9,0.9,0.9}{
            \makecell[{{p{\linewidth}}}]{
                \texttt{\tiny{[GM$|$GM]}}
                {'word': 'боєць', 'assignment': 'team'}
            }
        }
    }
    & & \\ \\

    \theutterance \stepcounter{utterance}  
    & & & \multicolumn{2}{p{0.3\linewidth}}{
        \cellcolor[rgb]{0.9,0.9,0.9}{
            \makecell[{{p{\linewidth}}}]{
                \texttt{\tiny{[GM$|$GM]}}
                {'word': 'заряд', 'assignment': 'innocent'}
            }
        }
    }
    & & \\ \\

    \theutterance \stepcounter{utterance}  
    & & & \multicolumn{2}{p{0.3\linewidth}}{
        \cellcolor[rgb]{0.9,0.9,0.9}{
            \makecell[{{p{\linewidth}}}]{
                \texttt{\tiny{[GM$|$GM]}}
                \texttt{заряд} \\
            }
        }
    }
    & & \\ \\

    \theutterance \stepcounter{utterance}  
    & & & \multicolumn{2}{p{0.3\linewidth}}{
        \cellcolor[rgb]{0.9,0.9,0.9}{
            \makecell[{{p{\linewidth}}}]{
                \texttt{\tiny{[GM$|$GM]}}
                {'word': 'оливка', 'assignment': 'opponent'}
            }
        }
    }
    & & \\ \\

    \theutterance \stepcounter{utterance}  
    & & & \multicolumn{2}{p{0.3\linewidth}}{
        \cellcolor[rgb]{0.9,0.9,0.9}{
            \makecell[{{p{\linewidth}}}]{
                \texttt{\tiny{[GM$|$GM]}}
                {'word': "м'ята", 'assignment': 'opponent'}
            }
        }
    }
    & & \\ \\

    \theutterance \stepcounter{utterance}  
    & & & \multicolumn{2}{p{0.3\linewidth}}{
        \cellcolor[rgb]{0.9,0.9,0.9}{
            \makecell[{{p{\linewidth}}}]{
                \texttt{\tiny{[GM$|$GM]}}
                {'hidden': {'team': ['тарілка', 'страйк', 'лід', 'єгипет', 'вашингтон', 'вчитель', 'храм'], 'innocent': ['весна', 'тінь', 'готель', 'костюм', 'орел', 'ліс'], 'opponent': ['агент', 'токіо', 'центр'], 'assassin': ['бомба']}, 'revealed': {'team': {'team': ['струна', 'боєць'], 'innocent': ['заряд'], 'opponent': ['театр'], 'assassin': []}, 'opponent': {'team': [], 'innocent': [], 'opponent': ['бермуди', 'європа', 'оливка', "м'ята"], 'assassin': []}}}
            }
        }
    }
    & & \\ \\

    \theutterance \stepcounter{utterance}  
    & & \multicolumn{4}{p{0.6\linewidth}}{
        \cellcolor[rgb]{0.9,0.9,0.9}{
            \makecell[{{p{\linewidth}}}]{
                \texttt{\tiny{[P1$\langle$GM]}}
                \texttt{Слово боєць було відгадано правильно. Слово заряд було відгадано, але це нейтральне слово. Хід вашого товариша по команді закінчився.} \\
\texttt{Тепер надайте ще одну підказку, пов'язану з деякими з ваших командних слів, що залишилися, та список пов'язаних цільових слів. Пам'ятайте: почніть підказку з 'ПІДКАЗКА: ', перейдіть на новий рядок і почніть ваш список цільових слів, розділених комами, з 'ЦІЛІ: '. Зверніть увагу: деякі слова було вилучено зі списків порівняно з попередніми запитами.} \\
\\ 
\texttt{Ваші командні слова, що залишилися: тарілка, страйк, лід, єгипет, вашингтон, вчитель, храм.} \\
\texttt{Слова, що залишилися для вашого суперника: агент, оливка, токіо, центр, м'ята.} \\
\texttt{Слова{-}відволікачі, що залишилися: весна, тінь, готель, костюм, орел, ліс.} \\
\texttt{Слова{-}вбивці, що залишилися: бомба.} \\
            }
        }
    }
    & \\ \\

    \theutterance \stepcounter{utterance}  
    & \multicolumn{4}{p{0.6\linewidth}}{
        \cellcolor[rgb]{0.9,0.9,0.9}{
            \makecell[{{p{\linewidth}}}]{
                \texttt{\tiny{[P1$\rangle$GM]}}
                \texttt{ПІДКАЗКА: спорт} \\
\texttt{ЦІЛІ: страйк} \\
            }
        }
    }
    & & \\ \\

    \theutterance \stepcounter{utterance}  
    & & & \multicolumn{2}{p{0.3\linewidth}}{
        \cellcolor[rgb]{0.9,0.9,0.9}{
            \makecell[{{p{\linewidth}}}]{
                \texttt{\tiny{[GM$|$GM]}}
                \texttt{спорт} \\
            }
        }
    }
    & & \\ \\

    \theutterance \stepcounter{utterance}  
    & & & \multicolumn{2}{p{0.3\linewidth}}{
        \cellcolor[rgb]{0.9,0.9,0.9}{
            \makecell[{{p{\linewidth}}}]{
                \texttt{\tiny{[GM$|$GM]}}
                ['страйк']
            }
        }
    }
    & & \\ \\

    \theutterance \stepcounter{utterance}  
    & & & \multicolumn{2}{p{0.3\linewidth}}{
        \cellcolor[rgb]{0.9,0.9,0.9}{
            \makecell[{{p{\linewidth}}}]{
                \texttt{\tiny{[GM$|$GM]}}
                {'word': 'страйк', 'assignment': 'team'}
            }
        }
    }
    & & \\ \\

    \theutterance \stepcounter{utterance}  
    & & \multicolumn{4}{p{0.6\linewidth}}{
        \cellcolor[rgb]{0.9,0.9,0.9}{
            \makecell[{{p{\linewidth}}}]{
                \texttt{\tiny{[GM$\rangle$P2]}}
                \texttt{Слово боєць було відгадано правильно. Слово заряд було відгадано, але це нейтральне слово. Ваш хід закінчився.} \\
\texttt{Тепер надайте ще один список слів, розділених комами, щонайменше з 1 і щонайбільше з 1 слів із наведеного нижче списку слів, які найкраще пов'язані або найбільш тісно асоціюються зі словом 'спорт'. Пам'ятайте, що потрібно почати вашу відповідь з ВІДГАДКА: . Зауваження: деякі слова було вилучено зі списку порівняно з попередніми запитами.} \\
\\ 
\texttt{готель, храм, страйк, бомба, весна, орел, костюм, лід, ліс, єгипет, тарілка, вашингтон, вчитель, центр, тінь, токіо, агент} \\
            }
        }
    }
    & \\ \\

    \theutterance \stepcounter{utterance}  
    & & & \multicolumn{4}{p{0.6\linewidth}}{
        \cellcolor[rgb]{0.9,0.9,0.9}{
            \makecell[{{p{\linewidth}}}]{
                \texttt{\tiny{[GM$\langle$P2]}}
                \texttt{ВІДГАДКА: лід} \\
            }
        }
    }
     \\ \\

    \theutterance \stepcounter{utterance}  
    & & & \multicolumn{2}{p{0.3\linewidth}}{
        \cellcolor[rgb]{0.9,0.9,0.9}{
            \makecell[{{p{\linewidth}}}]{
                \texttt{\tiny{[GM$|$GM]}}
                ['лід']
            }
        }
    }
    & & \\ \\

    \theutterance \stepcounter{utterance}  
    & & & \multicolumn{2}{p{0.3\linewidth}}{
        \cellcolor[rgb]{0.9,0.9,0.9}{
            \makecell[{{p{\linewidth}}}]{
                \texttt{\tiny{[GM$|$GM]}}
                {'word': 'лід', 'assignment': 'team'}
            }
        }
    }
    & & \\ \\

    \theutterance \stepcounter{utterance}  
    & & & \multicolumn{2}{p{0.3\linewidth}}{
        \cellcolor[rgb]{0.9,0.9,0.9}{
            \makecell[{{p{\linewidth}}}]{
                \texttt{\tiny{[GM$|$GM]}}
                {'word': 'центр', 'assignment': 'opponent'}
            }
        }
    }
    & & \\ \\

    \theutterance \stepcounter{utterance}  
    & & & \multicolumn{2}{p{0.3\linewidth}}{
        \cellcolor[rgb]{0.9,0.9,0.9}{
            \makecell[{{p{\linewidth}}}]{
                \texttt{\tiny{[GM$|$GM]}}
                {'word': 'токіо', 'assignment': 'opponent'}
            }
        }
    }
    & & \\ \\

    \theutterance \stepcounter{utterance}  
    & & & \multicolumn{2}{p{0.3\linewidth}}{
        \cellcolor[rgb]{0.9,0.9,0.9}{
            \makecell[{{p{\linewidth}}}]{
                \texttt{\tiny{[GM$|$GM]}}
                {'hidden': {'team': ['тарілка', 'страйк', 'єгипет', 'вашингтон', 'вчитель', 'храм'], 'innocent': ['весна', 'тінь', 'готель', 'костюм', 'орел', 'ліс'], 'opponent': ['агент'], 'assassin': ['бомба']}, 'revealed': {'team': {'team': ['струна', 'боєць', 'лід'], 'innocent': ['заряд'], 'opponent': ['театр'], 'assassin': []}, 'opponent': {'team': [], 'innocent': [], 'opponent': ['бермуди', 'європа', 'оливка', "м'ята", 'центр', 'токіо'], 'assassin': []}}}
            }
        }
    }
    & & \\ \\

    \theutterance \stepcounter{utterance}  
    & & \multicolumn{4}{p{0.6\linewidth}}{
        \cellcolor[rgb]{0.9,0.9,0.9}{
            \makecell[{{p{\linewidth}}}]{
                \texttt{\tiny{[P1$\langle$GM]}}
                \texttt{Слово лід було відгадано правильно. Хід вашого товариша по команді закінчився.} \\
\texttt{Тепер надайте ще одну підказку, пов'язану з деякими з ваших командних слів, що залишилися, та список пов'язаних цільових слів. Пам'ятайте: почніть підказку з 'ПІДКАЗКА: ', перейдіть на новий рядок і почніть ваш список цільових слів, розділених комами, з 'ЦІЛІ: '. Зверніть увагу: деякі слова було вилучено зі списків порівняно з попередніми запитами.} \\
\\ 
\texttt{Ваші командні слова, що залишилися: тарілка, страйк, єгипет, вашингтон, вчитель, храм.} \\
\texttt{Слова, що залишилися для вашого суперника: агент, токіо, центр.} \\
\texttt{Слова{-}відволікачі, що залишилися: весна, тінь, готель, костюм, орел, ліс.} \\
\texttt{Слова{-}вбивці, що залишилися: бомба.} \\
            }
        }
    }
    & \\ \\

    \theutterance \stepcounter{utterance}  
    & \multicolumn{4}{p{0.6\linewidth}}{
        \cellcolor[rgb]{0.9,0.9,0.9}{
            \makecell[{{p{\linewidth}}}]{
                \texttt{\tiny{[P1$\rangle$GM]}}
                \texttt{ПІДКАЗКА: столиця} \\
\texttt{ЦІЛІ: вашингтон, єгипет} \\
            }
        }
    }
    & & \\ \\

    \theutterance \stepcounter{utterance}  
    & & & \multicolumn{2}{p{0.3\linewidth}}{
        \cellcolor[rgb]{0.9,0.9,0.9}{
            \makecell[{{p{\linewidth}}}]{
                \texttt{\tiny{[GM$|$GM]}}
                \texttt{столиця} \\
            }
        }
    }
    & & \\ \\

    \theutterance \stepcounter{utterance}  
    & & & \multicolumn{2}{p{0.3\linewidth}}{
        \cellcolor[rgb]{0.9,0.9,0.9}{
            \makecell[{{p{\linewidth}}}]{
                \texttt{\tiny{[GM$|$GM]}}
                ['вашингтон', 'єгипет']
            }
        }
    }
    & & \\ \\

    \theutterance \stepcounter{utterance}  
    & & & \multicolumn{2}{p{0.3\linewidth}}{
        \cellcolor[rgb]{0.9,0.9,0.9}{
            \makecell[{{p{\linewidth}}}]{
                \texttt{\tiny{[GM$|$GM]}}
                {'word': 'вашингтон', 'assignment': 'team'}
            }
        }
    }
    & & \\ \\

    \theutterance \stepcounter{utterance}  
    & & & \multicolumn{2}{p{0.3\linewidth}}{
        \cellcolor[rgb]{0.9,0.9,0.9}{
            \makecell[{{p{\linewidth}}}]{
                \texttt{\tiny{[GM$|$GM]}}
                {'word': 'єгипет', 'assignment': 'team'}
            }
        }
    }
    & & \\ \\

    \theutterance \stepcounter{utterance}  
    & & \multicolumn{4}{p{0.6\linewidth}}{
        \cellcolor[rgb]{0.9,0.9,0.9}{
            \makecell[{{p{\linewidth}}}]{
                \texttt{\tiny{[GM$\rangle$P2]}}
                \texttt{Слово лід було відгадано правильно. Ваш хід закінчився.} \\
\texttt{Тепер надайте ще один список слів, розділених комами, щонайменше з 1 і щонайбільше з 2 слів із наведеного нижче списку слів, які найкраще пов'язані або найбільш тісно асоціюються зі словом 'столиця'. Пам'ятайте, що потрібно почати вашу відповідь з ВІДГАДКА: . Зауваження: деякі слова було вилучено зі списку порівняно з попередніми запитами.} \\
\\ 
\texttt{готель, храм, страйк, бомба, весна, орел, костюм, ліс, єгипет, тарілка, вашингтон, вчитель, тінь, агент} \\
            }
        }
    }
    & \\ \\

    \theutterance \stepcounter{utterance}  
    & & & \multicolumn{4}{p{0.6\linewidth}}{
        \cellcolor[rgb]{0.9,0.9,0.9}{
            \makecell[{{p{\linewidth}}}]{
                \texttt{\tiny{[GM$\langle$P2]}}
                \texttt{ВІДГАДКА: вашингтон, єгипет} \\
            }
        }
    }
     \\ \\

    \theutterance \stepcounter{utterance}  
    & & & \multicolumn{2}{p{0.3\linewidth}}{
        \cellcolor[rgb]{0.9,0.9,0.9}{
            \makecell[{{p{\linewidth}}}]{
                \texttt{\tiny{[GM$|$GM]}}
                ['вашингтон', 'єгипет']
            }
        }
    }
    & & \\ \\

    \theutterance \stepcounter{utterance}  
    & & & \multicolumn{2}{p{0.3\linewidth}}{
        \cellcolor[rgb]{0.9,0.9,0.9}{
            \makecell[{{p{\linewidth}}}]{
                \texttt{\tiny{[GM$|$GM]}}
                {'word': 'вашингтон', 'assignment': 'team'}
            }
        }
    }
    & & \\ \\

    \theutterance \stepcounter{utterance}  
    & & & \multicolumn{2}{p{0.3\linewidth}}{
        \cellcolor[rgb]{0.9,0.9,0.9}{
            \makecell[{{p{\linewidth}}}]{
                \texttt{\tiny{[GM$|$GM]}}
                {'word': 'вашингтон', 'assignment': 'team'}
            }
        }
    }
    & & \\ \\

    \theutterance \stepcounter{utterance}  
    & & & \multicolumn{2}{p{0.3\linewidth}}{
        \cellcolor[rgb]{0.9,0.9,0.9}{
            \makecell[{{p{\linewidth}}}]{
                \texttt{\tiny{[GM$|$GM]}}
                {'word': 'єгипет', 'assignment': 'team'}
            }
        }
    }
    & & \\ \\

    \theutterance \stepcounter{utterance}  
    & & & \multicolumn{2}{p{0.3\linewidth}}{
        \cellcolor[rgb]{0.9,0.9,0.9}{
            \makecell[{{p{\linewidth}}}]{
                \texttt{\tiny{[GM$|$GM]}}
                {'word': 'єгипет', 'assignment': 'team'}
            }
        }
    }
    & & \\ \\

    \theutterance \stepcounter{utterance}  
    & & & \multicolumn{2}{p{0.3\linewidth}}{
        \cellcolor[rgb]{0.9,0.9,0.9}{
            \makecell[{{p{\linewidth}}}]{
                \texttt{\tiny{[GM$|$GM]}}
                {'word': 'агент', 'assignment': 'opponent'}
            }
        }
    }
    & & \\ \\

    \theutterance \stepcounter{utterance}  
    & & & \multicolumn{2}{p{0.3\linewidth}}{
        \cellcolor[rgb]{0.9,0.9,0.9}{
            \makecell[{{p{\linewidth}}}]{
                \texttt{\tiny{[GM$|$GM]}}
                {'hidden': {'team': ['тарілка', 'страйк', 'вчитель', 'храм'], 'innocent': ['весна', 'тінь', 'готель', 'костюм', 'орел', 'ліс'], 'opponent': [], 'assassin': ['бомба']}, 'revealed': {'team': {'team': ['струна', 'боєць', 'лід', 'вашингтон', 'єгипет'], 'innocent': ['заряд'], 'opponent': ['театр'], 'assassin': []}, 'opponent': {'team': [], 'innocent': [], 'opponent': ['бермуди', 'європа', 'оливка', "м'ята", 'центр', 'токіо', 'агент'], 'assassin': []}}}
            }
        }
    }
    & & \\ \\

    \theutterance \stepcounter{utterance}  
    & & \multicolumn{4}{p{0.6\linewidth}}{
        \cellcolor[rgb]{0.9,0.9,0.9}{
            \makecell[{{p{\linewidth}}}]{
                \texttt{\tiny{[P1$\langle$GM]}}
                \texttt{Слова вашингтон, єгипет було відгадано правильно. Хід вашого товариша по команді закінчився.} \\
\texttt{Тепер надайте ще одну підказку, пов'язану з деякими з ваших командних слів, що залишилися, та список пов'язаних цільових слів. Пам'ятайте: почніть підказку з 'ПІДКАЗКА: ', перейдіть на новий рядок і почніть ваш список цільових слів, розділених комами, з 'ЦІЛІ: '. Зверніть увагу: деякі слова було вилучено зі списків порівняно з попередніми запитами.} \\
\\ 
\texttt{Ваші командні слова, що залишилися: тарілка, страйк, вчитель, храм.} \\
\texttt{Слова, що залишилися для вашого суперника: агент.} \\
\texttt{Слова{-}відволікачі, що залишилися: весна, тінь, готель, костюм, орел, ліс.} \\
\texttt{Слова{-}вбивці, що залишилися: бомба.} \\
            }
        }
    }
    & \\ \\

    \theutterance \stepcounter{utterance}  
    & \multicolumn{4}{p{0.6\linewidth}}{
        \cellcolor[rgb]{0.9,0.9,0.9}{
            \makecell[{{p{\linewidth}}}]{
                \texttt{\tiny{[P1$\rangle$GM]}}
                \texttt{ПІДКАЗКА: освіта} \\
\texttt{ЦІЛІ: вчитель} \\
            }
        }
    }
    & & \\ \\

    \theutterance \stepcounter{utterance}  
    & & & \multicolumn{2}{p{0.3\linewidth}}{
        \cellcolor[rgb]{0.9,0.9,0.9}{
            \makecell[{{p{\linewidth}}}]{
                \texttt{\tiny{[GM$|$GM]}}
                \texttt{освіта} \\
            }
        }
    }
    & & \\ \\

    \theutterance \stepcounter{utterance}  
    & & & \multicolumn{2}{p{0.3\linewidth}}{
        \cellcolor[rgb]{0.9,0.9,0.9}{
            \makecell[{{p{\linewidth}}}]{
                \texttt{\tiny{[GM$|$GM]}}
                ['вчитель']
            }
        }
    }
    & & \\ \\

    \theutterance \stepcounter{utterance}  
    & & & \multicolumn{2}{p{0.3\linewidth}}{
        \cellcolor[rgb]{0.9,0.9,0.9}{
            \makecell[{{p{\linewidth}}}]{
                \texttt{\tiny{[GM$|$GM]}}
                {'word': 'вчитель', 'assignment': 'team'}
            }
        }
    }
    & & \\ \\

    \theutterance \stepcounter{utterance}  
    & & & \multicolumn{2}{p{0.3\linewidth}}{
        \cellcolor[rgb]{0.9,0.9,0.9}{
            \makecell[{{p{\linewidth}}}]{
                \texttt{\tiny{[GM$|$GM]}}
                \texttt{opponent has won} \\
            }
        }
    }
    & & \\ \\

\end{supertabular}
}

\end{document}
